\centerline{\bf \Huge ДЗ. Обратный ход. Перебор. Чётность}
\centerline{\bf \Huge }

\begin{flushright}
        \scriptsize{\textbf{Сайгаки знают математику?!}}
\end{flushright}

\problem{\textbf{1}} Сайгак загадал интересное число разделил его на 4, прибавил 17, разделил на 5, прибавил 4, умножил на 11, разделил на 2, отнял 3 и у него получилось 52! Какое интересное число загадал сайгак?

\problem{\textbf{2}} АА и ученик А играли на чокопайки. Сначал АА проиграл половину своих чокопаек и отдал их ученику, затем ученик проиграл половину своих чокопаек и отдал их АА и так далее. В итоге АА проиграл 52 раза и ученик А проиграл 52 раза. У кого осталось больше чокопаек, если изначально у АА и у ученика А было одинаковое количество чокопаек?

\problem{\textbf{3}} Сколько пятизначных чисел может составить НеУченик из цифр 0, 2 и 5, если цифры могут повторяться?

\problem{\textbf{4}} Сколько чётных чисел от 13 до 2025 имеют различные цифры?

\problem{\textbf{5}} У АА есть огромная копилка, в которую он закинул какое-то количество пятитысячных купюр. Ученики закинули в эту же копилку стодолларовые купюры, которых было на 52 больше, чем пятитысячных купюр АА. Может ли в копилке оказаться ровно 2025 купюр?

\problem{\textbf{6}} В 201 кабинете ЭТЛ сидят 52 ученика. АА может выбрать любых трёх учеников и попросить их "встать-сесть" (тот, кто сидит должен встать, тот, кто стоит должен сесть). Может ли АА сделать так, чтобы все 52 ученика встали?